\emph{Origin and numbering.} Clauses (i) to (iii) are clause (i) of Theorem synthesis.light-ray-good-part-of-a-spinor-is-two-chirality-arrays, split in three (its projectors, its flux identity, and the two complex channels, whose count it gave in its clause (iii)); clause (iv) is clause (i) of Theorem synthesis.slot-module-light-ray-arrays-are-weyl-components. The arguments below are moved from those proofs; the step that Lorentz-covariant mass terms act from the right, argued there, is proved here.

\emph{Setting.} Every statement is about linear maps and quadratic forms on the real vector space $\Cl^{+}(1,3)$ with basis $1$, $\gamma_{0}\gamma_{i}$, $\gamma_{i}\gamma_{j}$ ($i < j$), $I_{4}$, or on $\Cl(1,3)$ with its sixteen blades. The script builds the rational matrices of the maps below and decides each identity and each dimension by exact rank computation. A rotor carrying $\gamma_{0} \pm \gamma_{3}$ to another pair of null directions conjugates every map below to the corresponding map and preserves ranks, so the choice of $k$ and $m$ is no loss.

\emph{(i).} $B^{2} = \gamma_{3}\gamma_{0}\gamma_{3}\gamma_{0} = -\gamma_{3}^{2}\gamma_{0}^{2} = 1$, so $L_{B}^{2} = 1$ and $\Cl^{+}(1,3)$ is the direct sum of the $\pm1$ eigenspaces of $L_{B}$. Directly, $k\gamma_{0} = 1 + B$, $m\gamma_{0} = 1 - B$, $km = 2 + 2B$ and $mk = 2 - 2B$, so $\Lambda(k) = \tfrac12 k\gamma_{0} = \tfrac12(1 + B) = \tfrac14 km$ and $\Lambda(m) = \tfrac12 m\gamma_{0} = \tfrac12(1 - B) = \tfrac14 mk$ (Theorem \claimref{field-theory.second-quantisation}(1) at $a_{0} = 1$). From $B^{2} = 1$ they are idempotents with sum $1$ and product $\tfrac14(1 - B^{2}) = 0$ in either order. $\tilde{B} = \gamma_{0}\gamma_{3} = -B$, so the reverse of $\Lambda(m)$ is $\tfrac12(1 + B)$, and $\gamma_{0}B\gamma_{0} = \gamma_{0}\gamma_{3} = -B$, so $\gamma_{0}\tfrac12(1 + B)\gamma_{0} = \tfrac12(1 - B) = \Lambda(m)$; likewise for $\Lambda(k)$. If $B\psi = -\psi$ then $\psi = \tfrac12 m\gamma_{0}\psi$ and $m\psi = \tfrac12 m^{2}\gamma_{0}\psi = 0$; conversely $m\psi = 0$ gives $\Lambda(k)\psi = \tfrac14 km\psi = 0$. The script finds $\ker L_{m}$ and the $(-1)$-eigenspace of $L_{B}$ both of dimension $4$ with intersection of dimension $4$, and $L_{m}L_{\Lambda(m)} = 0$, $(L_{B} + 1)L_{\Lambda(m)} = 0$, $\mathrm{rank}\,L_{\Lambda(m)} = 4$; the same with $k$ and the $(+1)$-eigenspace.

\emph{(ii).} If $\psi = \Lambda(k)\chi = \tfrac12 k\gamma_{0}\chi$, then $J_{\psi} = \psi\gamma_{0}\tilde{\psi} = \tfrac14 k(\gamma_{0}\chi\gamma_{0}\tilde{\chi}\gamma_{0})k$, and for any vector $v$, $kvk = 2\ip{k}{v}k - vk^{2} = 2\ip{k}{v}k$, so $J_{\psi}$ is a multiple of $k$ and $J_{\psi}\cdot k = 0$: the image of $\Lambda(k)$ carries no flux through $N_{k}$. The script computes the full Gram matrix of $\psi \mapsto J_{\psi}\cdot k = \langle\psi\gamma_{0}\tilde{\psi}k\rangle_{0}$ and finds it equal to $2\Lambda(m)^{T}\Lambda(m)$ exactly, so $J_{\psi}\cdot k = 2|\Lambda(m)\psi|^{2}$, positive definite on the image of $\Lambda(m)$ (rank $4$) and zero on that of $\Lambda(k)$; the flux through $N_{m}$ is $2|\Lambda(k)\psi|^{2}$, likewise, and as a contrast the flux through a hyperplane with normal $\gamma_{0}$ is not carried by either half alone. On the spinor $\psi = 1 + 2\gamma_{1}\gamma_{2} + 3\gamma_{0}\gamma_{3}$ it finds $J_{\psi}\cdot k = 20 = 2|\Lambda(m)\psi|^{2}$. The canonical part of a spinor field on $N_{k}$ is therefore the part that carries its current through $N_{k}$, $\Lambda(m)\psi$. The spinor $\Lambda(k)$ of a plane wave of momentum $k$ (Theorem \claimref{field-theory.second-quantisation}(1)) is constant along the generators of $N_{k}$, and its canonical data live on $N_{m}$. (This is the good component of light-front quantisation, Kogut and Soper, cited from memory as orientation.)

\emph{(iii).} $I_{4}^{2} = -1$ and $(I_{4}\sigma_{3})^{2} = -1$. $I_{4}$ commutes with every bivector, so $L_{I_{4}}$ commutes with $L_{B}$; right and left multiplications commute, so $R_{I_{4}\sigma_{3}}$ commutes with $L_{B}$. Each therefore commutes with $L_{\Lambda(m)} = \tfrac12(1 - L_{B})$ and preserves its image $V_{\Lambda}$, which, of real dimension $4$, is a complex vector space of dimension $2$ for the complex structure given by either unit (the script checks both squares, both commutations and $U L_{\Lambda(m)} = L_{\Lambda(m)} U L_{\Lambda(m)}$ for both units $U$, and, as a contrast, that left multiplication by $\gamma_{1}\gamma_{3}$ does not commute with $L_{\Lambda(m)}$). The same holds for the image of $\Lambda(k)$.

\emph{(iv).} With $\ip{k}{m} = 2$ and $Q(k) = Q(m) = 0$ the reciprocal vectors of $k$ and $m$ in the plane of $k$ and $m$ are $\tfrac12 m$ and $\tfrac12 k$, so every vector is $a = \tfrac12 m\,\ip{k}{a} + \tfrac12 k\,\ip{m}{a} + a_{\perp}$ with $a_{\perp}$ orthogonal to both, and correspondingly $\nabla = \tfrac12 m\,(k\cdot\nabla) + \tfrac12 k\,(m\cdot\nabla) + \nabla_{\perp}$. The derivative $k\cdot\nabla$ is along the generators, inside $N_{k}$, and $\nabla_{\perp}$ is along $N_{k}$ as well; only $m\cdot\nabla$ leaves the plane. Since $k^{2} = Q(k) = 0$, $k\nabla = \tfrac12 km\,(k\cdot\nabla) + k\nabla_{\perp} = 2\Lambda(k)(k\cdot\nabla) + k\nabla_{\perp}$ has no derivative off $N_{k}$. By (i) $k\Lambda(k) = \tfrac12 k^{2}\gamma_{0} = 0$, and $k$ anticommutes with $\gamma_{1}$ and $\gamma_{2}$ because $\ip{k}{\gamma_{j}} = 0$; so for every multivector $X$, $k\psi X = k(\Lambda(k) + \Lambda(m))\psi X = k\Lambda(m)\psi X$.

\emph{Covariant maps are right multiplications.} Let $f : \Cl^{+}(1,3) \to \Cl(1,3)$ be linear and commute with left multiplication by every even element, $f(Y\psi) = Yf(\psi)$. Then $f(\psi) = f(\psi\cdot 1) = \psi f(1)$, so $f$ is right multiplication by $f(1)$, and conversely every right multiplication commutes with every left one. Lorentz covariance of a mass term is commutation with $\psi \mapsto R\psi$ for every rotor $R$, and the rotors span $\Cl^{+}(1,3)$ ($e^{\theta X}$ for the basis bivectors $X$ give $1$ and the six bivectors, and products of two commuting ones give $I_{4}$), so a Lorentz-covariant mass term is a right multiplication on the spinor index, together with linear maps on any internal index: $\psi \mapsto \psi\gamma_{0}$ for a Dirac mass (Theorem \claimref{field-theory.second-quantisation}), and $\psi \mapsto \psi\sigma_{1}$, the charge conjugate, composed with a symmetric matrix on the internal index, for a Majorana mass. The script confirms this by solving the linear equations $FL_{X} = L_{X}F$ for an unknown $16 \times 8$ matrix $F$, $X$ the six basis bivectors (which generate $\Cl^{+}(1,3)$ as an algebra): the solution space has dimension $16$ and is spanned by the sixteen right multiplications by blades.

\emph{The constraint.} Write a massive spinor equation as $\nabla\psi I_{4}\sigma_{3} = \mathcal{M}(\psi)$ with $\mathcal{M}$ such a right multiplier. Multiplying on the left by $k$ gives
\[
2\Lambda(k)\,(k\cdot\nabla)\psi\,I_{4}\sigma_{3} + k\nabla_{\perp}\psi\,I_{4}\sigma_{3} = k\,\mathcal{M}(\psi).
\]
Here $\Lambda(k)$ commutes with $k\cdot\nabla$ and with right multiplication, $k\nabla_{\perp}\Lambda(k)\psi = -\nabla_{\perp}k\Lambda(k)\psi = 0$, and $k\,\mathcal{M}(\Lambda(k)\psi) = 0$ for every right multiplier. So the equation reads $2(k\cdot\nabla)(\Lambda(k)\psi)I_{4}\sigma_{3} = -k\nabla_{\perp}\Lambda(m)\psi\,I_{4}\sigma_{3} + k\,\mathcal{M}(\Lambda(m)\psi)$: along each generator it fixes $\Lambda(k)\psi$ from $\Lambda(m)\psi$ (up to its value at one end of the generator), and the mass enters only here. The flux identity (ii), which defines the canonical part, involves no field equation. So the free data on $N_{k}$ are $\Lambda(m)\psi$ for every mass matrix. (The mass independence of the null-plane data is the familiar property of light-front quantisation, cited from memory as orientation; the argument above does not use it.) The script checks $k^{2} = 0$, $km = 2 + 2B$, the splitting of each $\gamma_{\mu}$, the anticommutation of $k$ with $\gamma_{1}, \gamma_{2}$, $\mathrm{rank}(L_{k}L_{\Lambda(k)}) = 0$, $\mathrm{rank}(L_{k}L_{\Lambda(m)}) = 4$, and $\mathrm{rank}(L_{k}R_{X}L_{\Lambda(k)}) = 0$ for each of the sixteen basis blades $X$, and, as a contrast, that the left multiplier $\gamma_{3}$ does not pass: $k\gamma_{3}\Lambda(k) = -2\Lambda(k)$.

\emph{What is and is not derived.} All four clauses are exact identities in $\Cl(1,3)$ resting only on the field theory's projector $\Lambda(a)$ and Dirac mass term. They do not compute the anticommutator of the restricted field, which the results that use them compute in their own settings.
